\documentclass[11pt]{article}
\usepackage[T1]{fontenc}\usepackage{lmodern,amsmath,amssymb,amsthm,microtype}
\usepackage[margin=1in]{geometry}\usepackage[hidelinks]{hyperref}
\setlength\emergencystretch{1.5em}
\newtheorem{theorem}{Theorem}
\newtheorem{proposition}[theorem]{Proposition}
\title{Omitted branches in a paired Jacobsthal lower sieve\\
and finite divisor tests on translated intervals}
\author{\small \href{mailto:admin@cosmosgeek.org}{admin@cosmosgeek.org}\qquad \href{mailto:b321014@iiit-bh.ac.in}{b321014@iiit-bh.ac.in}}
\date{Public working paper for specialist review}
\begin{document}\maketitle
\begin{abstract}
Let $L=\log z$. Sieving $n,n+2$ on $(z,z^2-2]$ through $z/L^A$ leaves $O_A(z^2\log L/L^3)$ composite exceptions in addition to ordinary twin pairs. We examine the signed lower tree obtained by applying the admission rule of the OpenAI Jacobsthal preprint to two forbidden classes. Under the compact prime-path occupation hypothesis stated below, we prove that the omitted-node count exceeds every fixed multiple of $z^2/\log^2 z$ for all sufficiently large $z$. This obstruction persists under exact stops confined to a fixed exponent window.
Separately, an exact divisor expansion gives the parity of the twin count
on any interval of even midpoints. For every finite cutoff in this
expansion we construct arbitrarily high intervals on which the truncated
count is odd and every candidate endpoint is composite. This statement
uses elementary congruences alone.
\end{abstract}
\section{Composite exceptions}
Write $P^-(m)$ for the least prime factor of $m>1$, $L=\log z$, $I_z=(z,z^2-2]\cap\mathbb Z$ and $H=|I_z|$. Fix $A>0$ and put $z_0=z/L^A$. Define
\[
 S(y)=\#\{n\in I_z:P^-(n(n+2))>y\},\qquad T(z)=\#\{n\in I_z:n,n+2\text{ prime}\}.
\]
Both endpoints exceed $z$ and are at most $z^2$, hence $S(z)=T(z)$.
\begin{theorem}For every fixed $A>0$,
\[
 0\le S(z_0)-T(z)\ll_A H\log L/L^3.
\]
\end{theorem}
\begin{proof}
Since $z_0^3>z^2$ eventually, every surviving composite endpoint is $rs$, with $z_0<r\le s$ prime and $r\le z$. For fixed $r$, apply the two-linear-form upper-bound sieve to $s,rs\pm2$ on $[1,z^2/r]$ through a sufficiently small fixed power of $z$. The length is at least $z$. Every sieving prime other than $2$ has two distinct roots since $r>z_0$; at $2$ there is one. Each root class has interval counting error $O(1)$. A sufficiently small fixed power as divisor level absorbs even the sum of absolute divisor-weighted boundary errors. The standard Selberg upper bound gives $O(z^2/(rL^2))$ uniformly. Partial summation and the prime number theorem with its classical remainder give $\sum_{z_0<r\le z,\ r\ \rm prime}1/r\ll_A(\log L)/L$. Sum over $r$ and allow either endpoint to be composite; squares are included.
\end{proof}
\section{The signed tree}
Let $w=L/(\log L)^2$, $B=\log_w z_0$, $a_*=1/100$. Nodes carry decreasing primes in $(w,z_0]$, each labelled by $p\mid n$ or $p\mid n+2$. For their product $d$ put $r_d=\log_w(H/d)+2-a_*$. The root and the other even-depth nodes admit every available smaller prime. An odd-depth node with product $d'$ admits a new prime of exponent $u=\log_w p$ exactly when
\[
                     r_{d'}-u\ge\max(2u,u+2).                 \tag{1}
\]
The initial upper endpoint is included. Let $N_\nu$ count the integers satisfying the node congruences and $P^-(n(n+2))>w$. Let $S_\nu$ additionally exclude all available smaller prime hits. Partition by the least smaller prime hit. The two labels at an odd prime are disjoint. Thus $S_\nu=N_\nu-\sum_{\nu'} S_{\nu'}$ over all possible children. Expanding to the first rejected nodes gives
\[
             S(z_0)=\mathcal L+\Delta,\qquad
 \mathcal L=\sum_{\nu\ {\rm admitted}}(-1)^{\mathrm{depth}(\nu)}N_\nu,
 \quad\Delta=\sum_{\nu\ {\rm first\ rejected}}S_\nu.           \tag{2}
\]
Every first rejected child has even depth. Node counts in $\Delta$ carry their multiplicities.

\paragraph{Occupation input.} Retain the single colour $p\mid n$, exponents $>2$, and the standard rule $r_{d'}-u\ge2u$ at odd-depth parents. For a nonempty node put $s_d=r_d/\log_w P^-(d)$. We use the following consequence of the compact-window prime-occupation proposition in \cite[\S5, Proposition 5.5(iv)]{OAI}. For each fixed $b\ge20$,
\[
 \liminf_{z\to\infty}B^2
 \sum_{\substack{d'\ {\rm odd\ depth}\\2.5b<r_{d'}<2.8b\\2.05<s_{d'}<2.2}}
 \frac1{d'}\ge c_0 b^2                                      \tag{3}
\]
with $c_0>0$ independent of $b$. Indeed $r_1/B\to2$ and the odd-state limit on the rectangle is a fixed positive multiple of $r(1-s^{-2})\,dr\,ds$; the support restriction $s<r/2$ holds there. These paths retain precisely the same $1/d'$ weights as in (1) since the two admission rules agree above exponent $2$. The following theorem states the dependence on (3) explicitly.

\begin{theorem}\label{main}Assuming (3), $\Delta/(H/L^2)\to+\infty$ and $\mathcal L/(H/L^2)\to-\infty$. The latter assertion also holds if exact survivor counts replace the lower-tree expansions at a prefix-free set of nonroot stops whose last exponents are bounded above by a fixed constant.
\end{theorem}
\begin{proof}
To each parent in (3) append a prime $p=w^u$ of the same colour with $b<u<1.1b$. The preceding exponent is $>2.5b/2.2>1.13b$, but $r_{d'}-u<1.8b<2u$, so the child is first rejected. For $d=d'p$ put $v=\log_w(H/d)$. Direct substitution gives
\[
 v=r_{d'}-u-2+a_*,\qquad v-u>0.3b-2,\quad v<2u,\quad uv<1.98b^2. \tag{4}
\]
For fixed $b$, each $d$ exceeds $z$ eventually. In its count use only $n=dk$ with $k>p$ prime. The first endpoint then survives all primes $<p$. We claim that, summed over these $d$, the number of witnesses with $P^-(dk+2)\ge p$ is
\[
 \sum_d W(p)\#\{k>p:k\text{ prime},\ dk\in I_z\}+o_b(H/L^2),
 \qquad W(p)=\prod_{2<\ell<p}(1-1/(\ell-1)).                  \tag{5}
\]
First discard products whose largest prime is at most $R_*=\exp(L^{1/3})$. Their reciprocal mass is bounded by
\[
 z^{-L^{-1/3}}\prod_{\ell\le R_*,\ \ell\ {\rm prime}}
       (1+\ell^{-1+L^{-1/3}})\ll e^{-L^{2/3}} L^{O(1)} .
\]
Each numerical $d$ determines its decreasing tuple; it contributes $O(H/d)$ integers since $d<H$. This removal is negligible on both sides of (5).

For the remaining products let $P$ be the largest prime, and $d=Pa$. For fixed $a,k$ and $R<P\le2R$, the admissible $P$ lie in an interval $J\subseteq[R,2R]$. Here $p=P^-(a)$, and every later prefix gap is affine in $\log P$. The admission tests, rectangle (3), decreasing order and $Pak\in I_z$ each restrict $\log P$ to an interval. The largest prime cannot be the appended $p$, or the last prime of the parent once $P>R_*$. For each odd $\ell<p$ the class $-2(ak)^{-1}\pmod\ell$ is forbidden. Uniformly for these classes and intervals, and for fixed $D>0$,
\[
 \#\{P\in J:P\text{ prime},\,P^-(Pak+2)\ge p\}
 =W(p)\#\{P\in J:P\text{ prime}\}+O_{D,b}(RL^{-D}).           \tag{6}
\]
Indeed truncate inclusion--exclusion at adjacent even and odd orders $J_0,J_0+1$, with $J_0\asymp C_D\log L$. The principal discrepancy is bounded by the reciprocal factorial tail with parameter $\sum_{2<\ell<p}1/(\ell-1)=O_b(\log\log L+1)$, hence by $O_{D,b}(L^{-D-2})$. Each squarefree modulus occurs at most once in each sum. All moduli are at most $p^{J_0+1}=\exp(O_{D,b}((\log L)^2))$. As $R\gg\exp(L^{1/3})$, maximal Bombieri--Vinogradov and partial summation give the error in (6) with any required logarithmic precision, including on short $J$ as an absolute error. The prime number theorem allows the unrestricted prime count on the right. There are $O((H/R)L)$ possible pairs $a,k$ per dyadic block and $O(L)$ blocks. Choosing $D>4$ proves (5).

By (4) and the prime number theorem the inner prime count in (5) is $(1+o_b(1))H/(dv\log w)$ uniformly. Mertens' theorem gives $W(p)\sim2C_2e^{-\gamma}/(u\log w)$, $C_2=\prod_{\ell>2,\ \ell\ {\rm prime}}(1-1/(\ell-1)^2)$. Also $\sum_{w^b<p<w^{1.1b}}1/p\to\log(1.1)$. Applying (3) and (4) in (5), using $B\log w\sim L$, gives
\[
 \liminf_{z\to\infty}\frac{\Delta_b}{H/L^2}
 \ge \frac{2C_2e^{-\gamma} c_0}{1.98}\log(1.1)>0
\]
for the selected omissions in this band. The lower bound is independent of $b$. Take $M$ disjoint bands $b_j=20\cdot2^j$, keeping $M$ fixed while taking the lower limit. Nonnegativity in (2) permits their contributions to be added, and $M$ is arbitrary. The ordinary upper-bound sieve through a small fixed power of $z$ yields $S(z_0)\ll H/L^2$, proving both limits. For stops with bounded terminal exponent take all $b_j$ above that bound plus one. All ancestors of each selected rejected child remain above the stop window, so its omission persists after exact substitution.
\end{proof}

\setcounter{equation}{6}
\section{An exact parity formula on translated intervals}
The preceding estimate weighs the branches omitted by a particular
lower sieve on the anchored interval $I_z$. A different approach to
the twin count groups reciprocal divisor conditions and works modulo
two. We derive the expansion and determine exactly what every fixed
cutoff can miss. The positions of the intervals in the counterexample
below are free to vary.

For integers $2\le U\le V$ define
\[
 \mathcal T(U,V)=\#\{U\le j\le V:2j-1,\,2j+1\text{ are prime}\}.
\]
Let $\mathbf p(n)$ be the primality indicator, $\omega(n)$ the number of
distinct prime divisors, and $\mu$ the M\"obius function. One derivation
starts with the square roots of one modulo $m^2-1$ for $m=2j$.
An odd prime divides at most one of $u-1,u+1$. Thus $u^2=1$
modulo an odd prime power has exactly two solutions, the two signs.
The Chinese remainder theorem gives $2^{\omega(m-1)+\omega(m+1)}$
roots modulo $m^2-1$, since the neighbors are odd and coprime.
The four distinct multipliers $1,-1,m,-m$ form a subgroup of the
roots and act freely. Signs may be exchanged independently on the two
endpoints, since $m\equiv1\bmod(m-1)$ and $m\equiv-1\bmod(m+1)$.
The orbit count is odd precisely when both neighbors are prime powers.
Multiplication by the two squarefree indicators then selects ordinary
primes. For a different root-based criterion using the product of
invertible quadratic residues, see Aebi--Cairns \cite[Theorem 2]{AC}.

For a squarefree odd $n>1$, there are $2^{\omega(n)-1}$ unordered
factor pairs, with one member smaller than $\sqrt n$.
Their number is odd exactly when $n$ is prime. Hence
\begin{equation}\label{eq:divpar}
 \mathcal T(U,V)\equiv\sum_{j=U}^V\mu^2(2j-1)\mu^2(2j+1)
 \sum_{\substack{a\mid 2j-1\\a^2<2j}}1
 \sum_{\substack{c\mid 2j+1\\c^2<2j}}1 \pmod2 .
\end{equation}
The strict bound $2j$ works on both sides. For nonsquare endpoints it
selects precisely the smaller divisors, since an odd square cannot
equal $2j$. A square endpoint is killed by its squarefree indicator.

Expand $\mu^2(n)=\sum_{e^2\mid n}\mu(e)$. Modulo two the only terms
are squarefree $e$, each with coefficient one. Call an ordered choice
$(a,e;c,f)$ a row, with $a,c$ odd positive and $e,f$ odd squarefree
positive; all four parameters may be restricted to
$\lfloor\sqrt{2V+1}\rfloor$. Define
\[
 A=\operatorname{lcm}(a,e^2),\quad C=\operatorname{lcm}(c,f^2),
 \quad P=AC,\quad K=\max\{U,(\max(a^2,c^2)+1)/2\}.
\]
The row counts $K\le j\le V$ with $A\mid 2j-1$ and $C\mid2j+1$.
If $\gcd(A,C)>1$ the count is zero, since an odd common prime would
divide two. Otherwise the count is in one residue $r\bmod P$:
\begin{equation}\label{eq:rowcrt}
   2r\equiv1\pmod A,\qquad 2r\equiv-1\pmod C.
\end{equation}
Swapping $(a,e)$ with $(c,f)$ preserves $P,K$ and sends $r$ to
$-r$. The only fixed compatible row has $a=c=e=f=1$. For every
other compatible row $P>1$ is odd and $\gcd(r,P)=1$, so its two
orientations have distinct residues.

For $1\le r<P$ define
\[
 B_P(t;r)=\mathbf1_{\{\min(r,P-r)\le (t\bmod P)<\max(r,P-r)\}},
\]
using the least nonnegative remainder. Every complete period contains
both residues once. In a partial period from $0$ to $t\bmod P$ the
number of hits is odd exactly when $B_P(t;r)=1$. Thus a reciprocal
pair of rows contributes $B_P(V;r)+B_P(K-1;r)\bmod2$ when $K\le V$.
Distinct choices having the same $A,C$ retain their multiplicities.

\begin{proposition}\label{prop:mixed}
Let $\mathcal M(U,V)$ be the sum modulo two of these boundary terms,
once per unordered compatible reciprocal pair with $A,C>1$ and
$K\le V$. Then
\[
 \mathcal T(U,V)\equiv V-U+1+
 \mathbf p(2U-1)+\mathbf p(2V+1)+\mathcal M(U,V)\pmod2 .
\]
\end{proposition}
\begin{proof}
The trivial row contributes $V-U+1$. At a fixed $j$, the full
single-endpoint expansion in \eqref{eq:divpar} equals the primality
indicator modulo two. Removing its trivial term leaves
$\mathbf p(2j-1)+1$ on the left and $\mathbf p(2j+1)+1$ on the right.
These give all the one-sided rows. On summing, each interior primality
indicator occurs twice, leaving exactly the displayed endpoint terms.
The reciprocal pairing evaluates the remaining mixed rows.
\end{proof}

\section{Intervals on which the finite divisor test fails}
For an integer $R\ge1$ let $\mathcal T_R(U,V)$ be the sum of the
row counts in the expanded right side of \eqref{eq:divpar} modulo
two, restricted to $a,c,e,f\le R$. The squarefree inversion is part
of this truncation. Set
\[
       L_R=\operatorname{lcm}\{i^2:1\le i\le R,\ i\text{ odd}\}.
\]
\begin{theorem}\label{th:period}
For every $U>R^2$ with $U\ge2$ one has
$\mathcal T_R(U,U+L_R-1)\equiv1\bmod2$. For every $b\bmod L_R$ there
are arbitrarily high pairwise disjoint intervals of this length,
with $U\equiv b\bmod L_R$, on which all $L_R+1$ possible odd
endpoints have at least two distinct prime factors. On each such
interval the full twin count is zero and $\mathcal M=1\bmod2$.
\end{theorem}
\begin{proof}
Every retained row has $K=U$. Both $A$ and $C$ divide $L_R$,
because each allowed $a,e^2,c,f^2$ divides it. Their coprimality
on a compatible row gives $P=AC\mid L_R$. A nontrivial reciprocal
pair therefore has $2L_R/P$ hits, an even number. The trivial row
contributes the odd length $L_R$. This proves the first assertion.

For each $0\le i\le L_R$ choose two odd primes $p_i,q_i>R$,
all $2(L_R+1)$ primes distinct. Prescribe
\[
    2U\equiv1-2i\pmod{p_iq_i}\quad(0\le i\le L_R),\qquad
    U\equiv b\pmod{L_R}.
\]
All the moduli are pairwise coprime. The Chinese remainder theorem
gives an unbounded progression of solutions. Choose $U>R^2$ in it
and pass successively to starts beyond the preceding interval.
Every endpoint $2(U+i)-1$ is divisible by $p_iq_i$, proving the
second assertion. Both outer primality indicators in
Proposition~\ref{prop:mixed} vanish, so its formula gives
$\mathcal M=1\bmod2$ on these intervals.
\end{proof}

For example $R=3$ gives $L_R=9$. The interval $566\le j\le574$
has ten possible odd endpoints from $1131$ through $1149$. In
increasing order they have the proper divisors
\[
 3,\ 11,\ 5,\ 3,\ 17,\ 7,\ 3,\ 5,\ 31,\ 3 .
\]
Thus $\mathcal T=0$ on this interval whereas $\mathcal T_3=1\bmod2$.

The omitted rows can reverse the truncated answer even after every
retained modulus has completed a whole number of periods. The
conclusion also holds along a sequence of increasing cutoffs by
choosing each new interval beyond the last. The construction concerns
translated intervals. Prescribing a relation between their height
and the cutoff would require a separate argument. For each \emph{fixed}
$[U,V]$, taking $R\ge\lfloor\sqrt{2V+1}\rfloor$ recovers the
full expansion in \eqref{eq:divpar}. No prime-path
occupation estimate enters Theorem~\ref{th:period}.

\paragraph{Provenance.} The original Twin Prime Research chats supplied
the paired-tree adaptation, the squarefree reciprocal expansion and the
translated-interval construction. Codex checked and edited the proofs.
Theorem~\ref{main} addresses the specified signed lower tree, and
Theorem~\ref{th:period} addresses a separate finite congruence test.
The existence of twin pairs at arbitrarily large cutoffs and the
priority of these specializations remain open in this work.
The original modular calculation is preserved in the
\href{https://github.com/cosmosgeek0/Twin-Prime/blob/main/history/hosted-proofs-supplement.md}{project source supplement, part 2}.
\begin{thebibliography}{9}
\bibitem{OAI}OpenAI, \emph{A quadratic bound for Jacobsthal's function} (2026). \href{https://github.com/openai/math/blob/main/preprints/A-quadratic-bound-for-Jacobsthals-function-September-25-2026/build/sections/prime-paths.tex}{Prime-path occupation proposition}.
\bibitem{IK}H. Iwaniec and E. Kowalski, \emph{Analytic Number Theory}, AMS, 2004.
\bibitem{AC}C. Aebi and G. Cairns, \emph{A property of twin primes},
Integers 12 (2012), A7.
\href{https://math.colgate.edu/~integers/m7/m7.pdf}{Published article}.
\bibitem{record}\href{https://github.com/cosmosgeek0/Twin-Prime/blob/main/history/hosted-research-responses-2026-11-25.md}{Original project research and independent check}, section on the paired Jacobsthal lower sieve.
\end{thebibliography}
\end{document}
